\begin{name}
    {\tenchude}
    {\tendethi}
    {\tentruong}
    {\thoigian}
\end{name}
\Opensolutionfile{ansbook}[ans/ansbook-5]
\TN
\Opensolutionfile{ans}[ans/ans-5-T]
\begin{ex}%Câu 1
Họ nguyên hàm của hàm số $ f(x)=\dfrac{1}{2x+3}$ là
\choice
{$ 3\ln \left| 2x+3\right|+C$}
{$\dfrac{1}{3}\ln \left| 2x+3\right|+C$}
{$ 2\ln \left| 2x+3\right|+C$}
{\True $\dfrac{1}{2} \ln \left| 2x+3\right|+C$}
\loigiai{
Ta có $\displaystyle\int{f(x)}\mathrm{\,d}x=\dfrac{1}{2}\ln \left| 2x+3\right|+C$.
}
\end{ex}

\begin{ex}%Câu 2
Diện tích $ S$ của hình phẳng giới hạn bởi đồ thị hàm số $ y=f(x)$, liên tục trên $\left[a;b\right]$ trục hoành và hai đường thẳng $ x=a,\,x=b$ ($ a<b$) cho bởi công thức:
\choice
{\True $ S=\displaystyle\int\limits_a^b{\left| f(x)\right|dx}$}
{$ S=\displaystyle\int\limits_a^b{f(x)dx}$}
{$ S=\pi\displaystyle\int\limits_a^b{\left| f(x)\right|dx}$}
{$ S=\pi\displaystyle\int\limits_a^b{f^2(x)dx}$}
\loigiai{
Diện tích $ S$ của hình phẳng là: $ S=\displaystyle\int\limits_a^b{\left| f(x)\right|dx}$}
\end{ex}

\begin{ex}%Câu 3
Sau khi kiểm tra sức khỏe tổng quát, kết quả số cân nặng của học sinh lớp 12A sĩ số 40 HS được thể hiện trong bảng số liệu sau: (Đơn vị: kg)
\begin{center}
\begin{tabular}{|c|c|c|c|c|c|}
    \hline
    Cân nặng & $[40; 50)$ & $[50; 60)$ & $[60; 70)$ & $[70; 80)$ & $[80; 90)$ \\
    \hline
    Số học sinh & 7 & 12 & 12 & 7 & 2 \\
    \hline
\end{tabular}
\end{center}
Tứ phân vị thứ nhất của mẫu số liệu trên gần nhất với giá trị nào trong các giá trị sau?
\choice
{$ 50$}
{$ 50,5$}
{\True $ 52,5$}
{$ 55,5$}
\loigiai{
Tứ phân vị thứ nhất của dãy số liệu thuộc nhóm $\left[50;60\right)$ nên tứ phân vị thứ nhất của mẫu số liệu là $Q_1=50+\dfrac{\dfrac{40}{4}-7}{12}\left(60-50\right)=52,5$}
\end{ex}

\begin{ex}%Câu 4
Trong không gian $Oxyz$ , phương trình của đường thẳng đi qua điểm $ A\left(1;2;\,-1\right)$ và có vectơ chỉ phương $\vec{u}=\left(1;\,3;\,2\right)\,$ là
\choice
{$\dfrac{x+1}{1}\,=\,\dfrac{y\,+\,3}{2}\,=\,\dfrac{z\,+\,2}{-1}$}
{$\dfrac{x-1}{1}\,=\,\dfrac{y\,-\,3}{2}\,=\,\dfrac{z\,-\,2}{-1}$}
{$\dfrac{x+1}{1}\,=\,\dfrac{y\,+\,2}{3}\,=\,\dfrac{z\,-\,1}{2}$}
{\True $\dfrac{x-1}{1}\,=\,\dfrac{y\,-\,2}{3}\,=\,\dfrac{z\,+\,1}{2}$}
\loigiai{
Đường thẳng đi qua điểm $ A\left(1;2;\,-1\right)$ và có vectơ chỉ phương $\vec{u}=\left(1;\,3;\,2\right)\,$ có phương trình là: $\dfrac{x-1}{1}\,=\,\dfrac{y\,-\,2}{3}\,=\,\dfrac{z\,+\,1}{2}$.}
\end{ex}

\begin{ex}%Câu 5
Hàm số số $ y=\dfrac{ax+b}{cx+d}$ có đồ thị như hình bên dưới:\\
\centerline{
    \begin{tikzpicture}[line join=round, line cap=round,>=stealth,thick]
\tikzset{every node/.style={scale=0.9}}
\draw[->] (-3.1,0)--(3.1,0) node[below left] {$x$};
\draw[->] (0,-2.1)--(0,4.1) node[below left] {$y$};
\draw (0,0) node [below left] {$O$}
(-1,0)node[below left]{$-1$}
(0,2)node[above right]{$2$};
\draw[dashed,thin] (-0.99,-2)--(-0.99,4);
\begin{scope}
\clip (-3,-2) rectangle (3,4);
\draw[samples=200,domain=-4:-1.01,smooth,variable=\x] plot (\x,{(2*(\x)+1)/(1*(\x)+1)});
\draw[samples=200,domain=-0.99:4,smooth,variable=\x] plot (\x,{(2*(\x)+1)/(1*(\x)+1)});
\draw[dashed,thin] (-4,2/1)--(4,2/1);
\end{scope}
\end{tikzpicture}
}
Đường tiệm cận đứng của đồ thị là đường thẳng có phương trình
\choice
{$ x=1$}
{$ x=2$}
{$ x=-2$}
{\True $ x=-1$}
\loigiai{
Dựa vào đồ thị ta có đường tiệm cận đứng của đồ thị là đường thẳng $ x=-1$.}
\end{ex}

\begin{ex}%Câu 6
Tập nghiệm của bất phương trình $\log_{\dfrac{1}{2}}\left(x-2\right)\le 1$ là:
\choice
{\True $\left[\dfrac{5}{2};\,\,+\infty\right)$}
{$\left(\dfrac{5}{2};\,\,+\infty\right)$}
{$\left(-\infty ;{\log_2}5\right)$}
{$\left(-\infty ;\dfrac{5}{2}\right)$}
\loigiai{
Điều kiện: $ x-2>0\Leftrightarrow x>2$.\\
Bất phương trình: $\log_{\dfrac{1}{2}}\left(x-2\right)\le 1\Leftrightarrow x-2\ge\dfrac{1}{2}\Leftrightarrow x\ge\dfrac{5}{2}$.\\
Kết hợp với điều kiện ta có tập nghiệm $ T=\left[\dfrac{5}{2};\,\,+\infty\right)$.}
\end{ex}

\begin{ex}%Câu 7
Trong không gian $Oxyz$ , mặt phẳng nào sau đây nhận $\overrightarrow{n}=\left(1;2;3\right)$ làm vectơ pháp tuyến?
\choice
{$x+2y+3=0$}
{\True $x+2y+3z=0$}
{$y+2z+3=0$}
{$x+2z+3=0$}
\loigiai{
Mặt phẳng $x+2y+3z=0$ có vectơ pháp tuyến là $\overrightarrow{n}=\left(1;2;3\right)$ .}
\end{ex}

\begin{ex}%Câu 8
Cho hình chóp $ S.ABCD$ có đáy $ ABCD$ là hình chữ nhật tâm $ I$ và cạnh bên $ SA$ vuông góc với đáy. Khẳng định nào sau đây đúng?
\choice
{\True $\left(SCD\right)\perp\left(SAD\right)$}
{$\left(SBC\right)\perp\left(SIA\right)$}
{$\left(SDC\right)\perp\left(SAI\right)$}
{$\left(SBD\right)\perp\left(SAC\right)$}
\loigiai{
% {\color{red}HÌNH Ở ĐÂY}\\
Ta có: $CD\perp AD$ (Vì $ ABCD$ là hình chữ nhật) và $ SA\perp\left(ABCD\right)\Rightarrow SA\perp CD$\\
Mặt khác: $ SA\cap AD=A$ và $ SA,AD\subset\left(SAD\right)$$\Rightarrow CD\perp\left(SAD\right)$\\
Mà $CD\subset\left(SCD\right)$ nên $\left(SCD\right)\perp\left(SAD\right)$ .}
\end{ex}

\begin{ex}%Câu 9
Nghiệm của phương trình $\left(\dfrac{1}{5}\right)^{x^2-2x-3}=5^{x+1}$ là
\choice
{\True $ x=-1;\,x=2$}
{Vô nghiệm}
{$ x=1;\,x=2$}
{$ x=1;\,x=-2$}
\loigiai{
Phương trình đã cho tương đương $5^{-x^2+2x+3}=5^{x+1}\Leftrightarrow-x^2+x+2=0\Leftrightarrow\left[\begin{aligned}
& x=-1\\ 
& x=2.\\ 
\end{aligned}\right.$\\
Vậy phương trình có nghiệm $ x=-1;\,x=2$.}
\end{ex}

\begin{ex}%Câu 10
Cho cấp số cộng $\left(u_n\right)$ có $u_1=2$, $u_2=6$. Công sai của cấp số cộng bằng
\choice
{$8$}
{$-4$}
{3}
{\True 4}
\loigiai{
Vì $u_2=6\Leftrightarrow{u_1}+d=6\Leftrightarrow 2+d=6\Leftrightarrow d=4$ .}
\end{ex}

\begin{ex}%Câu 11
Cho hình lăng trụ tam giác $ ABC.A'B'C'$. Đặt $\overrightarrow{AA'}=\overrightarrow{a},\overrightarrow{AB}=\overrightarrow{b},\overrightarrow{AC}=\overrightarrow{c}$. Khi đó biểu diễn $\overrightarrow{BC'}$ theo các vectơ $\overrightarrow{a},\overrightarrow{b},\overrightarrow{c}$\\
\centerline{
\begin{tikzpicture}[scale=1, font=\footnotesize, line join=round, line cap=round, >=stealth]
    \def\ac{4} % cạnh AC
    \def\ab{2} % cạnh AB
    \def\h{4} % chiều cao
    \def\gocA{50} % góc A của đáy
    \path
    (0,0) coordinate (A)
    (\ac,0) coordinate (C)
    (-\gocA:\ab) coordinate (B)
    ($(A)+(80:\h)$) coordinate (A')
    ($(B)-(A)+(A')$) coordinate (B')
    ($(C)-(A)+(A')$) coordinate (C');
    \draw (A')--(A)--(B)--(C)--(C')--(A')--(B')--(C') (B)--(B');
    \draw[dashed] (A)--(C);
    \foreach \x/\g in {A/180,B/-90,C/0,A'/180,C'/0,B'/-40}\fill[red] (\x) circle (1pt)+(\g:3mm) node[black]{$ \x $};
\end{tikzpicture}
}
\choice
{$\overrightarrow{BC'}=-\overrightarrow{a}+\overrightarrow{b}+\overrightarrow{c}$}
{\True $\overrightarrow{BC'}=\overrightarrow{a}-\overrightarrow{b}+\overrightarrow{c}$}
{$\overrightarrow{BC'}=\overrightarrow{a}+\overrightarrow{b}+\overrightarrow{c}$}
{$\overrightarrow{BC'}=\overrightarrow{a}+\overrightarrow{b}-\overrightarrow{c}$}
\loigiai{
Do $ ABC.A'B'C'$ là hình lăng trụ nên $\overrightarrow{A'C'}=\overrightarrow{AC}$ nên ta có:\\
$\overrightarrow{BC'}=\overrightarrow{AC'}-\overrightarrow{AB}=\overrightarrow{AA'}+\overrightarrow{A'C'}-\overrightarrow{AB}=\overrightarrow{a}-\overrightarrow{b}+\overrightarrow{c}$}
\end{ex}

\begin{ex}%Câu 12
Cho hàm số $ y=f(x)$ có bảng biến thiên như sau :\\
\centerline{
    \begin{tikzpicture}
\tkzTabInit[espcl=2.5,lgt=1.5,nocadre=false]
{$x$/0.7,$f'(x)$/0.7,$f(x)$/2.1}
{$-\infty$,$0$,$2$,$+\infty$}
\tkzTabLine{,+,0,-,0,+,}
\tkzTabVar{-/$-\infty$,+/$4$,-/$0$,+/$+\infty$}
\end{tikzpicture}
}
Hàm số đã cho đồng biến trên khoảng nào dưới đây?
\choice
{$\left(-4\,;\,1\right)$}
{$\left(0\,;+\infty\right)$}
{\True $\left(-\infty ;\,0\,\right)$}
{$\left(0\,;2\right)$}
\loigiai{
Hàm số đã cho đồng biến trên khoảng $\left(-\infty ;\,0\,\right)$.}
\end{ex}
\Closesolutionfile{ans}
\TNTF
\Opensolutionfile{ans}[ans/ans-5-TF]

\begin{ex}%Câu 13
Cho hàm số $f(x)=\cos 2x+2x+1$ . 
\choiceTF
{\True $ f\left(\dfrac{\pi}{2}\right)=\pi $}
{Đạo hàm của hàm số đã cho là $f'(x)=2\sin 2x+2$}
{\True Nghiệm của phương trình $f'(x)=0$ trên đoạn $\left[-\dfrac{\pi}{2};\pi\right]$ là $ x=\dfrac{\pi}{4}$}
{Tổng giá trị lớn nhất và giá trị nhỏ nhất của hàm số trên đoạn $\left[-\dfrac{\pi}{2};\pi\right]$ bằng $ 2\pi $}
\loigiai{

a) Đúng: $ f\left(\dfrac{\pi}{2}\right)=\cos\left(2.\dfrac{\pi}{2}\right)+2.\dfrac{\pi}{2}+1=\pi $\\
b) Sai: Đạo hàm $f'(x)=-2\sin 2x+2$\\
c) Đúng: Phương trình $f'(x)=0\Leftrightarrow-2\sin 2x+2=0\Leftrightarrow\sin 2x=1\Leftrightarrow x=\dfrac{\pi}{4}+k\pi ,\,k\in\mathbb{Z}$\\
Vì $x\in\left[-\dfrac{\pi}{2};\pi\right]$ nên $x=\dfrac{\pi}{4}$\\
d) Sai: Ta có $f\left(-\dfrac{\pi}{2}\right)=-\pi $ ; $f\left(\dfrac{\pi}{4}\right)=\dfrac{\pi}{2}+1$ ; $f\left(\pi\right)=2\pi+2$ .\\
Vậy $\underset{\left[-\dfrac{\pi}{2};\,\pi\right]}{\max}\,f(x)=f\left(-\dfrac{\pi}{2}\right)=-\pi $ và $\underset{\left[-\dfrac{\pi}{2};\pi\right]}{\min}\,f(x)=f\left(\pi\right)=2\pi+2$\\
Khi đó tổng giá trị lớn nhất và nhỏ nhất của hàm số trên đoạn $\left[-\dfrac{\pi}{2};\pi\right]$ bằng $\pi+2$.}
\end{ex}

\begin{ex}%Câu 14
Các nhà kinh tế sử dụng đường cong Lorenz để minh họa sự phân phối thu nhập trong một quốc gia. Gọi $x$ là đại diện cho phần trăm số gia đình trong một quốc gia và $y$ là phần trăm tổng thu nhập, mô hình $y=x$ sẽ đại diện cho một quốc gia mà các gia đình có thu nhập như nhau. Đường cong Lorenz $y=f(x)$ , biểu thị sự phân phối thu nhập thực tế. Diện tích giữa hai mô hình này, với $0\le x\le 100$ , biểu thị “sự bất bình đẳng về thu nhập” của một quốc gia. Năm $2005$ , đường cong Lorenz của Hoa Kỳ có thể được mô hình hóa bởi hàm số:
$$y=\left(0,00061x^2+0,0218x+1,723\right)^2,0\le x\le 100$$
Trong đó $x$ được tính từ các gia đình nghèo nhất đến giàu có nhất\\
\textit{(Theo R.Larson, Brief Calculus: An Applied Approach, 8th edition, Cengage Learning, 2009)}
\choiceTF
{\True Tính theo thứ tự từ các gia đình nghèo nhất đến giàu nhất, tổng thu nhập thực tế của $60\%$ các gia đình đầu tiên chiếm chưa đến $30\%$ so với tổng thu nhập của toàn bộ các gia đình}
{Nếu sắp xếp các gia đình theo thứ tự từ nghèo nhất đến giàu nhất, rồi chia thành $10$ nhóm bằng nhau từ $1$ đến $10$ , tổng thu nhập của các gia đình trong nhóm $3$ chiếm khoảng $8,56\%$ tổng thu nhập của toàn bộ các gia đình}
{Sự bất bình đẳng về thu nhập của Hoa Kì năm $2005$ được xác định bởi công thức:\\
$\displaystyle\int\limits_0^{100}{\left[x-\left(0,00061x^2+0,0218x+1,723\right)^2\right]\mathrm{\,d}x}$}
{\True Sự bất bình đẳng về thu nhập của Hoa Kỳ năm $2005$ đã vượt quá $2000$}
\loigiai{

a) Đúng: Tính theo thứ tự từ các gia đình nghèo nhất đến giàu nhất, tổng thu nhập của $60\%$ các gia đình của đầu tiên chiếm tỷ lệ trong tổng thu nhập là: $f\left(60\right)=27,321529(\%)$ .\\
b) Sai: Nếu sắp xếp các gia đình theo thứ tự từ nghèo nhất đến giàu nhất, rồi chia thành $10$ nhóm bằng nhau từ $1$ đến $10$ , tổng thu nhập của các gia đình trong nhóm $3$ chiếm khoảng $8,56\%$ tổng thu nhập của toàn bộ các gia đình.\\
Nếu sắp xếp các gia đình theo thứ tự từ nghèo đến giàu, rồi chia thành $10$ nhóm bằng nhau, mỗi nhóm chiếm $10\%$ số gia đình của Hoa Kỳ.\\
Tổng thu nhập của $30\%$ số gia đình (là các gia đình thuộc nhóm $1,2,3$) chiếm tỷ lệ trong tổng thu nhập của tất cả các gia đình là: $f\left(30\right)=8,561476\,\,\left(\%\right)$ .\\
Tổng thu nhập của $20\%$ số gia đình (là các gia đình thuộc nhóm $1,2$) chiếm tỷ lệ trong tổng thu nhập của tất cả các gia đình là: $f\left(20\right)=5,774409\,\,\,\left(\%\right)$ .\\
Tỷ lệ của tổng thu nhập các gia đình nhóm thứ $3$ so với toàn bộ các gia đình là:\\
$f\left(30\right)-f\left(20\right)=2,787067(\%)$ .\\
c) Sai: Sự bất bình đẳng về thu nhập của Hoa Kì năm $2005$ được xác định bởi công thức:\\
$\displaystyle\int\limits_0^{100}{\left[x-\left(0,00061x^2+0,0218x+1,723\right)^2\right]\mathrm{\,d}x}$ .\\
Sự bất bình đẳng về thu nhập của Hoa Kì vào năm $2005$ là diện tích hình phẳng $S$ giới hạn bởi hai đồ thị:\\
$\left\{\begin{aligned}
& y=x\\ 
& y=\left(0,00061x^2+0,0218+1,723\right)^2\\ 
& x=0;x=100\\ 
\end{aligned}\right.$ $\Rightarrow S=\displaystyle\int\limits_0^{100}{\left|\left(0,00061x^2+0,0218x+1,723\right)^2-x\right|\mathrm{\,d}x}$ .\\
Cách 1:\\
Sử dụng máy tính cầm tay, ta thấy phương trình $\left(0,00061x^2+0,0218x+1,723\right)^2-x=0$ có hai ngiệm $x=a\,;x=b\,\,\left(a<b\right)$ thuộc $\left[0\,;100\right]$ .\\
Xét dấu biểu thức $g(x)=\left(0,00061x^2+0,0218x+1,723\right)^2-x$ ta suy ra:\\ $S=\displaystyle\int\limits_0^{100}{\left| g(x)\right|\mathrm{\,d}x}=\displaystyle\int\limits_0^a{\left| g(x)\right|\mathrm{\,d}x}+\displaystyle\int\limits_a^b{\left| g(x)\right|\mathrm{\,d}x}+\displaystyle\int\limits_b^{100}{\left| g(x)\right|\mathrm{\,d}x}$ .\\
$=\left|\displaystyle\int\limits_0^a{g(x)\mathrm{\,d}x}\right|+\left|\displaystyle\int\limits_a^b{g(x)\mathrm{\,d}x}\right|+\left|\displaystyle\int\limits_b^{100}{g(x)\mathrm{\,d}x}\right|$ $=\displaystyle\int\limits_0^a{g(x)\mathrm{\,d}x}-\displaystyle\int\limits_a^b{g(x)\mathrm{\,d}x}+\displaystyle\int\limits_b^{100}{g(x)\mathrm{\,d}x}$ .\\
Cách 2:\\
Sử dụng máy tính cầm tay ta được: $S=\displaystyle\int\limits_0^{100}{\left|\left(0,00061x^2+0,0218x+1,723\right)^2-x\right|\mathrm{\,d}x}\approx 2068,9$ .\\
Kiểm tra phép tính của đề bài, ta có: $\displaystyle\int\limits_0^{100}{\left[x-\left(0,00061x^2+0,0218x+1,723\right)^2\right]\mathrm{\,d}x}=2059,3131$ .\\
d) Đúng: Sự bất bình đẳng về thu nhập của Hoa Kỳ năm $2005$ đã vượt quá $2000$ .\\
Sự bất bình đẳng thu nhập của Hoa Kỳ năm $2005$ là:\\
$S=\displaystyle\int\limits_0^{100}{\left|\left(0,00061x^2+0,0218x+1,723\right)^2-x\right|\mathrm{\,d}x}\approx 2068,9$ .}
\end{ex}

\begin{ex}%Câu 15
Một công ty tham gia đấu thầu 2 dự án. Khả năng thắng thầu dự án 1 là 60\% và dự án 2 là 50\%. Khả năng thắng thầu cả hai dự án là 40\%. Gọi $ A$ và $ B$ lần lượt là biến cố công ty thắng thầu dự án 1 và dự án 2.
\choiceTF
{$ A$ và $ B$ là hai biến cố độc lập}
{\True Khả năng công ty thắng thầu đúng 1 dự án là 30\%}
{Xác suất công ty thắng thầu dự án 2 biết công ty đã thắng thầu dự án 1 là $\dfrac{1}{2}$}
{Xác suất công ty không thắng thầu dự án 2, biết công ty đã không thắng thầu dự án 1 là $\dfrac{1}{4}$}
\loigiai{
Ta có $P(A)=0,6\Rightarrow P\left(\bar{A}\right)=0,4;\quad P(B)=0,5\Rightarrow P\left(\bar{B}\right)=0,5$ và $P\left(AB\right)=0,4$ 

a) Sai: Vì $P\left(AB\right)\ne P(A).P(B)$ nên $A,B$ không độc lập .\\
b) Đúng: Do $ A\overline{B}$ và $\overline{A}B$ là hai biến cố xung khắc nên xác suất công ty thắng thầu đúng 1 dự án là: $P\left(A\bar{B}\right)+P\left(\bar{A}B\right)=P(A)-P\left(AB\right)+P(B)-P\left(AB\right)=0,6-0,4+0,5-0,4=0,3$\\
c) Sai: Xác suất công ty thắng thầu dự án 2 biết công ty đã thắng thầu dự án 1 là:\\
$P\left(B|A\right)=\dfrac{P\left(BA\right)}{P(A)}=\dfrac{0,4}{0,6}=\dfrac{2}{3}$ .\\
d) Sai: Xác suất công ty thắng thầu dự án 2, biết công ty đã không thắng thầu dự án 1 là:\\
$P\left(B|\bar{A}\right)=\dfrac{P\left(B\bar{A}\right)}{P\left(\bar{A}\right)}=\dfrac{P(B)-P\left(AB\right)}{P\left(\bar{A}\right)}=\dfrac{0,5-0,4}{0,4}=\dfrac{1}{4}$ .\\
Vậy xác suất công ty không thắng thầu dự án 2, biết công ty đã không thắng thầu dự án 1 là\\
$ P\left(\overline{B}\left|\overline{A}\right.\right)=1-P\left(B\left|\overline{A}\right.\right)=1-\dfrac{1}{4}=\dfrac{3}{4}$.}
\end{ex}

\begin{ex}%Câu 16
Trong không gian $Oxyz$ (đơn vị trên mỗi trục tính theo kilômét), một trạm thu phát sóng điện thoại di động được đặt ở vị trí $I\left(1;\,3;\,7\right)$ . Trạm thu phát sóng đó được thiết kế với bán kính phủ sóng là 3 km.\\
\centerline{\includegraphics[width=.4\textwidth]{images/5.16}}
\choiceTF
{Phương trình mặt cầu $(S)$ để mô tả ranh giới bên ngoài của vùng phù sóng trong không gian là $\left(x+1\right)^2+\left(y+3\right)^2+\left(z+7\right)^2=9$}
{\True Nếu người dùng điện thoại ở vị trí điểm $A\left(2;\,2;\,7\right)$ thì có thể sử dụng dịch vụ của trạm thu phát sóng đó}
{\True Nếu người dùng điện thoại ở vị trí có toạ độ $B\left(5;\,6;\,7\right)$ thì không thể sử dụng dịch vụ của trạm thu phát sóng đó}
{\True Tính theo đường chim bay, khoảng cách lớn nhất để một người ở vị trí có toạ độ $B\left(5;\,6;\,7\right)$ di chuyển được tới vùng phủ sóng theo đơn vị km là $8$ km}
\loigiai{
a) Sai: Phương trình mặt cầu $(S)$ tâm $I\left(1;\,3;\,7\right)$ bán kính 3km mô tả ranh giới bên ngoài của vùng phủ sóng trong không gian là $\left(x-1\right)^2+\left(y-3\right)^2+\left(z-7\right)^2=9$ .\\
b) Đúng: Ta có: $IA=\sqrt{\left(2-1\right)^2+\left(2-3\right)^2+\left(7-7\right)^2}=\sqrt{2}<3$ nên điểm $A$ nằm trong mặt cầu. Vì điểm $A$ nằm trong mặt cầu nên người dùng điện thoại ở vị trí có toạ độ $A\left(2;\,2;\,7\right)$ có thể sử dụng dịch vụ của trạm thu phát sóng đó.\\
c) Đúng: Ta có: $IB=\sqrt{\left(5-1\right)^2+\left(6-3\right)^2+\left(7-7\right)^2}=5>3$ nên điểm $B$ nằm ngoài mặt cầu. Vậy người dùng điện thoại ở vị trí có toạ độ $B\left(5;\,6;\,7\right)$ không thể sử dụng dịch vụ của trạm thu phát sóng đó.\\
d) Đúng: Ta có: $\overrightarrow{IB}\left(4;\,3;\,0\right);$ $IB=\sqrt{\left(5-1\right)^2+\left(6-3\right)^2+\left(7-7\right)^2}=5>3$ nên điểm $B$ nằm ngoài mặt cầu.\\
Phương trình đường thẳng $BI$ dạng: $\left\{\begin{aligned}
& x=1+4t\\ 
& y=3+3t\\ 
& z=7\\ 
\end{aligned}\right.$ .\\
Gọi mặt cầu $(S)\cap BI\equiv E$ suy ra tọa độ $E$ là nghiệm của hệ\\
$\left\{\begin{aligned}
& x=1+4t\\ 
& y=3+3t\\ 
& z=7\\ 
&{\left(x-1\right)^2}+\left(y-3\right)^2+\left(z-7\right)^2=9\\ 
\end{aligned}\right.\Leftrightarrow\left[\begin{aligned}
&\left\{\begin{aligned}
& t=\dfrac{3}{5}\\ 
& x=\dfrac{17}{5}\\ 
& y=\dfrac{24}{5}\\ 
& z=7\\ 
\end{aligned}\right.\Rightarrow E\left(\dfrac{17}{5};\,\dfrac{24}{5};7\right)\Rightarrow EB\approx 1,7\\ 
&\left\{\begin{aligned}
& t=-\dfrac{3}{5}\\ 
& x=-\dfrac{7}{5}\\ 
& y=\dfrac{6}{5}\\ 
& z=7\\ 
\end{aligned}\right.\Rightarrow E\left(-\dfrac{7}{5};\,\dfrac{6}{5};7\right)\Rightarrow EB=8\\ 
\end{aligned}\right.$\\
Vậy khoảng cách lớn nhất để một người ở vị trí có toạ độ $B\left(5;\,6;\,7\right)$ di chuyển được tới vùng phủ sóng theo đơn vị kilomet là $8$ km.}
\end{ex}
\Closesolutionfile{ans}
\TNSA
\Opensolutionfile{ans}[ans/ans-5-SA]
\begin{ex}%Câu 17
Cho hình hộp chữ nhật $ ABCD.A'B'C'D'$ có $ ABCD$ là hình vuông cạnh $ 2$ và góc giữa hai mặt phẳng $\left(AC'D'\right)$ và $\left(ABCD\right)$ bằng $30^{\text{o}}$. Tính khoảng cách giữa hai đường thẳng $ AD'$ và $A'B$. (Kết quả làm tròn đến chữ số thập phân thứ hai sau dấu phẩy)
\shortans{0,89}
\loigiai{
% {\color{red}HÌNH Ở ĐÂY}\\
Vì $\left(ABCD\right)\parallel\left(A'B'C'D'\right)$ nên góc giữa $\left(AC'D'\right)$ và mặt phẳng $\left(ABCD\right)$ bằng góc giữa $\left(AC'D'\right)$ và mặt phẳng $\left(A'B'C'D'\right)$.\\
Ta có $\heva{&\left(AC'D'\right)\cap\left(A'B'C'D'\right)=C'D'\\ &AD'\bot C'D' \\& A'D'\bot C'D'}$$\Rightarrow\left(\left(AC'D'\right);\left(A'B'C'D'\right)\right)=\left(AD';\,A'D'\right)=\widehat{AD'A'}=30^\circ $.\\
Trong tam giác vuông $ AA'D'$ có $ AA'=A'D'.\tan 30^\circ=\dfrac{2\sqrt{3}}{3}$.\\
Vì $A'B\parallel CD'$ nên $A'B\parallel\left(ACD'\right)$ suy ra $ d\left(A'B;\,AD'\right)=d\left(A'B;\,\left(ACD'\right)\right)=d\left(A';\,\left(ACD'\right)\right)$\\
Mặt khác $A'B\cap\left(ACD'\right)=I$ là trung điểm của $A'D$ nên $ d\left(A';\,\left(ACD'\right)\right)=d\left(D;\,\left(ACD'\right)\right)$.\\
Tại $ D$ ta có ba tia $ DA,\,DC,\,DD'$ đôi một vuông góc nên:\\
$\dfrac{1}{d^2\left(D;\,\left(ACD'\right)\right)}=\dfrac{1}{D{A^2}}+\dfrac{1}{D{C^2}}+\dfrac{1}{D{D'^2}}$$=\dfrac{1}{4}+\dfrac{1}{4}+\dfrac{3}{4}=\dfrac{5}{4}$\\
Suy ra $ d\left(D;\,\left(ACD'\right)\right)=\dfrac{2\sqrt{5}}{5}$ hay $ d\left(A'B;\,AD'\right)=\dfrac{2\sqrt{5}}{5}\approx 0,89$.}
\end{ex}

\begin{ex}%Câu 18
Công ty sản xuất "SmartTech" đang lên kế hoạch sản xuất hai loại sản phẩm A và B. Mỗi sản phẩm A đem lại lợi nhuận 800.000 đồng và mỗi sản phẩm B đem lại lợi nhuận 500.000 đồng. Tuy nhiên, việc sản xuất mỗi sản phẩm đòi hỏi nguyên vật liệu và công nhân khác nhau:\\
• Để sản xuất một sản phẩm A công ty cần sử dụng 2 kg nguyên vật liệu và 3 giờ lao động.\\
• Để sản xuất một sản phẩm B công ty cần sử dụng 1 kg nguyên vật liệu và 4 giờ lao động.\\
Hiện tại, công ty có sẵn 100 kg nguyên vật liệu và có thể sử dụng tối đa 180 giờ lao động. Công ty cần sản xuất $m$ sản phẩm A và $ n$ sản phẩm B để tối đa hóa lợi nhuận, đồng thời thỏa mãn các điều kiện về nguyên vật liệu và giờ lao động. Khi đó tổng $m+n$ bẳng bao nhiêu?
\shortans{56}
\loigiai{
Gọi $x$ là số lượng sản phẩm A và $y$ là số lượng sản phẩm B (Điều kiện: $x\ge 0;y\ge 0$)\\
Điều kiện về nguyên vật liệu: $2x+y\le 100$ .\\
Điều kiện về giờ lao động: $3x+4y\le 180$ .\\
Ta có hệ bất phương trình: $\left\{\begin{aligned}
& x\ge 0\\ 
& y\ge 0\\ 
& 2x+y\le 100\\ 
& 3x+4y\le 180\\ 
\end{aligned}\right.$\\
Miền nghiệm của hệ bất phương trình là miền không bị gạch sọc như hình vẽ dưới đây:\\
{\color{red}HÌNH Ở ĐÂY}\\
Ta có miền nghiệm của hệ bất phương trình là miền trong của tứ giác $OEFG$\\
Trong đó, $E\left(0;45\right),F(44;12),G\left(50;0\right)$ .\\
Tối ưu hóa lợi nhuận: Ta sẽ tính toán lợi nhuận $P=800.000x+500.000y$ tại các đỉnh của miền nghiệm ta được:\\
Tại $E\left(0;45\right)$ : $P=22.500.000$ ; $F\left(44;12\right)$ : $P=41.200.000$ ; $G\left(50;0\right)$ : $P=40.000.000$\\
Vậy để tối đa hóa lợi nhuận thì nhà máy cần sản xuất 44 sản phẩm A và 12 sản phẩm B.\\
Khi đó tổng số là $44+12=56$ sản phẩm hay $ m+n=56$.}
\end{ex}

\begin{ex}%Câu 19
Hệ thống định vị toàn cầu GPS (Global Positioning System) là một hệ thống cho phép xác định vị trí của một vật thể trong không gian. Trong cùng một thời điểm vị trí của một điểm $ M$ trong không gian sẽ được xác định bởi bốn vệ tinh cho trước nhờ các bộ thu phát tín hiệu đặt trên các vệ tinh. Giả sử trong không gian $ Oxyz$, tỉ lệ dài trên các trục là $10$ km tính cho một đơn vị tỉ lệ trên mỗi trục, cho bốn vệ tinh có tọa độ lần lượt là $ A\left(-1;2;-1\right),B\left(1;4;0\right),C\left(3;0;9\right),D\left(7;10;-1\right)$ đang tiến hành theo dõi vật thể $ M$ (coi là một chất điểm). Các vệ tinh dùng sóng điện từ có tần số 1MHz . Ở một thời điểm cả 4 vệ tinh bắn tín hiệu về $ M$ thì sau nhận được tín hiệu trả về ngay sau đó trong những khoảng thời gian là $t_A=0,4$ m/s, $t_B=0,2 $m/s, $ t_c=\dfrac{2}{3}$ m/s, $t_D=0,4$ m/s. Tính khoảng cách từ $ M$ đến $ O$ (biết rằng vận tốc của sóng điện từ bằng vận tốc ánh sáng $ c=3.10^8$ m/s và kết quả tính được làm tròn đến hàng phần trăm).
\shortans{6,78}
\loigiai{
Từ vệ tinh $ A$ sẽ phát tín hiệu đến điểm $ M$ và sẽ thu tín hiệu quay về $ A$\\
Từ đó ta có thể suy ra được: $ 2AM=c.t_A\Rightarrow AM=\dfrac{c.t_A}{2}$\\
Do đó $ AM=\dfrac{c.t_A}{2}=\dfrac{3.10^8.0,4.10^{-3}}{2}=60\,000$m$=60$km và ta quy đổi trong hệ trục $ Oxyz$ thì ta có $ AM=6$.\\
Tương tự: $ BM=3$, $ CM=10$, $ DM=6$.\\
Tại thời điểm bốn vị tinh bắn tín hiệu về điểm $ M$. Khi đó ta có hệ phương trình sau:\\
$\Leftrightarrow\left\{\begin{matrix}
{\left(a+1\right)^2}+\left(b-2\right)^2+\left(c+1\right)^2=36\,\,\,\,\,\,\,\,\,\,\,\,\,\,(1)\\
{\left(a-1\right)^2}+\left(b-4\right)^2+c^2=9\,\,\,\,\,\,\,\,\,\,\,\,\,\,\,\,\,\,\,\,\,\,\,\,\,\,\,\,\,(2)\\
{\left(a-3\right)^2}+b^2+\left(c-9\right)^2=100\,\,\,\,\,\,\,\,\,\,\,\,\,\,\,\,\,\,\,\,\,\,\,(3)\\
{\left(a-7\right)^2}+\left(b-10\right)^2+\left(c+1\right)^2=36\,\,\,\,\,\,\,\,\,\,\,(4)\\
\end{matrix}\right.\Leftrightarrow\left\{\begin{aligned}
&(2)-(1):-4a-4b-2c=-38\\ 
&(3)-(1):-8a+4b-20c=-20\\ 
&(4)-(1):-14a-14b=-144\\ 
\end{aligned}\right.$\\
Giải hệ phương trình này, ta tìm được $ a~=~3,b=6,c=1\Rightarrow M\left(3;\,6;\,1\right)$\\
Vậy khoảng cách từ $ M$ đến $ O$ là: $\sqrt{\left(3-0\right)^2+\left(6-0\right)^2+\left(1-0\right)^2}=\sqrt{46}\approx 6,78$.}
\end{ex}

\begin{ex}%Câu 20
Một biển quảng cáo có dạng hình vuông $ ABCD$ cạnh $ AB=4\text{m}$. Trên tấm biển đó có các đường tròn tâm $ A$ và đường tròn tâm $ B$ cùng bán kính $ R=4\text{m}$, hai đường tròn cắt nhau như hình vẽ. Chi phí để sơn phần gạch chéo là $ 150\, 000$ đồng/$\text{m}^2$, chi phí sơn phần màu đen là $ 100\,000$ đồng/ $\text{m}^2$ và chi phí để sơn phần còn lại là $ 250\,000$ đồng/$\text{m}^2$\\
\centerline{\includegraphics[width=.25\textwidth]{images/5.20}}\\
Hỏi số tiền để sơn biển quảng cáo theo cách trên (Đơn vị: triệu đồng và kết quả làm tròn đến chữ số thập phân thứ hai sau dấu phẩy)?
\shortans{2,2}
\loigiai{
Gọi $ I$ là giao điểm của 2 cung tròn $\overset\frown{AC};\overset\frown{BD}$ . Chọn gốc toạ độ $ A\left(0;0\right)$$\Rightarrow B\left(4,0\right)$\\
Xét cung tròn có phương trình $ y=\sqrt{16-x^2}$\\
Phần diện tích gạch chéo $ S=2.\displaystyle\int\limits_2^4\sqrt{16-x^2}\mathrm{\,d}x=16\left.\left(x+\dfrac{1}{2}\sin 2x\right)\right|_{\dfrac{\pi}{6}}^{\dfrac{\pi}{2}}=\dfrac{16\pi}{3}-4\sqrt{3}$\\
Phần diện tích màu đen: $ 2.\left(\dfrac{1}{4}\pi{4^2}-\dfrac{16\pi}{3}+4\sqrt{3}\right)=\dfrac{-8\pi}{3}+8\sqrt{3}$\\
Phần diện tích còn lại: $ 16-\left(\dfrac{16\pi}{3}-4\sqrt{3}+\dfrac{-8\pi}{3}+8\sqrt{3}\right)=16-\dfrac{8\pi}{3}-4\sqrt{3}$\\
Số tiền để sơn biển quảng cáo:\\
$\left(\dfrac{16\pi}{3}-4\sqrt{3}\right).150\text{000}\,\text{+}\left(\dfrac{-8\pi}{3}+8\sqrt{3}\right).100\text{000}+\left(16-\dfrac{8\pi}{3}-4\sqrt{3}\right).250\text{000}\,\approx 2,2$ triệu đồng.}
\end{ex}

\begin{ex}%Câu 21
Anh An thành lập một công ty sản xuất in ấn Sách Giáo Khoa chương trình "Chân trời sáng tạo". Nhằm tạo điều kiện cho các nhà sách tiêu thụ giá hợp lí, đơn giá mỗi bộ sách ban đầu được biểu diễn theo hàm $ p(x)=200-3x$ (nghìn đồng) với $ x$ là số lượng từng bộ sách bán ra và tổng chi phí sản xuất được biểu diễn theo hàm $ C(x)=75+\left(80+T\right)x-x^2$ (nghìn đồng) với mọi $ x$ thỏa $ 0\le x\le 40$, trong đó $ T$ (nghìn đồng) là mức thuế giá trị gia tăng VAT phải đóng trên từng số lượng bộ sách sản xuất ra mà công ty anh An phải chi trả. Xem như công ty anh An sản xuất đều đặn trong điều kiện lí tưởng, khi lợi nhuận của công ty đạt giá trị cao nhất thì tổng mức thuế phải chi trả cũng đồng thời cao nhất. Khi đó mức thuế của mỗi bộ sách mà công ty phải trả là bao nhiêu (đơn vị: nghìn đồng)?
\shortans{60}
\loigiai{
Trước hết ta có hàm chi phí sản xuất là:$ C(x)=75+\left(80+T\right)x-x^2$\\
Doanh thu của công ty anh An biểu diễn theo hàm $ R(x)=x.p(x)=x\left(200-3x\right)$\\
Lợi nhuận mà công ty anh An có được là:$ P(x)=R(x)-C(x)=-2x^2+\left(120-T\right)x-75$\\
Do cần xác định số lượng bộ sách bán ra đề lợi nhuận là cao nhất nên ta có:\\
$P'(x)=0\Leftrightarrow-4x^2+\left(120-T\right)=0\Leftrightarrow x=30-\dfrac{T}{4}$\\
Khi đó với thuế mỗi bộ sách là $ T$ thì tồng mức thuế công ty phải trả là $ G(T)=T\left(30-\dfrac{T}{4}\right)$\\
Khi lợi nhuận của công ty đạt giá trị cao nhất thì tổng mức thuế phải chi trả cũng đồng thời cao nhất nên ta suy ra $G'(T)=0\Leftrightarrow\dfrac{T}{2}-30=0\Leftrightarrow T=60$ (nghìn đồng).}
\end{ex}

\begin{ex}%Câu 22
Bạn Tuấn hằng ngày ăn sáng bằng xôi hoặc bún. Nếu hôm nay bạn ăn sáng bằng xôi thì xác suất để hôm sau bạn ăn sáng bằng bún là $ 0,7$. Xét một tuần mà thứ ba bạn ăn sáng bằng xôi. Biết xác suất để thứ năm tuần đó, bạn Tuấn ăn sáng bằng bún là $ 0,63$. Hỏi nếu hôm nay bạn ăn sáng bằng bún thì xác suất để hôm sau bạn ăn sáng bằng xôi là bao nhiêu?
\shortans{0,4}
\loigiai{
Giả sử nếu hôm nay bạn ăn sáng bằng bún thì xác suất để hôm sau bạn ăn sáng bằng xôi là $ x$ $\left(x<1\right)$.\\
Gọi $ A$ là biến cố “Thứ tư, bạn Tuấn ăn sáng bằng bún”,\\
$ B$ là biến cố “Thứ năm, bạn Tuấn ăn sáng bằng bún”, khi đó $ P(B)=0,63$\\
Ta có thứ ba bạn Tuấn ăn sáng bằng xôi nên $ P(A)=0,7$, $ P\left(\overline{A}\right)=1-0,7=0,3$\\
Vì nếu hôm nay bạn ăn sáng bằng bún thì xác suất để hôm sau bạn ăn sáng bằng xôi là $ x$ và ăn sáng bằng bún là $ 1-x$ hay $ P\left(B|A\right)=1-x$.\\
Ta có $ P\left(B|\overline{A}\right)=0,7$\\
Theo công thức xác suất toàn phần: $ P(B)=P(A).P\left(B|A\right)+P\left(\overline{A}\right).P\left(B|\overline{A}\right)$\\
$\Rightarrow 0,63=0,7.\left(1-x\right)+0,3.0,7$$\Rightarrow x=0,4$\\
Vậy nếu hôm nay bạn ăn sáng bằng bún thì xác suất để hôm sau bạn ăn sáng bằng xôi là $ 0,4$.
}
\end{ex}
\Closesolutionfile{ans}
\Closesolutionfile{ansbook}
\inputansbox{6,2,3}{ans/ans-4-T,ans/ans-4-TF,ans/ans-4-SA}